\section{Sparse Table}
\whenuse{Arreglo est\'atico y consultas idempotentes como m\'inimo, m\'aximo o
GCD. Elegir Disjoint Sparse Table si la operaci\'on no es idempotente.}

Para operaciones idempotentes como min, max o gcd. No requiere inversa y no
admite actualizaciones.

Esta adaptación del RMQ de KACTL (Johan Sannemo y pajenegod) conserva
\texttt{SparseTableMin(a)} y las consultas \textbf{inclusivas} \texttt{query(l,r)}
del handbook. \texttt{table[k][i]} guarda el mínimo de $2^k$ elementos desde
$i$; cada fila reserva sólo las posiciones donde ese bloque entra completo.

Para consultar, tomar $k=\lfloor\log_2(r-l+1)\rfloor$ y combinar el bloque que
empieza en $l$ con el que termina en $r$. Cubren todo el rango y pueden
superponerse: para mínimo, repetir elementos no cambia el resultado. Por eso
no sirve reemplazar mínimo por suma; en ese caso usar Disjoint Sparse Table.
El logaritmo se calcula con \texttt{\_\_builtin\_clz}, sin un arreglo auxiliar.

\complexity{build $O(n\log n)$; query $O(1)$; memoria $O(n\log n)$.}

\begin{lstlisting}[style=cpp]
struct SparseTableMin {
    vector<vector<ll>> table;

    explicit SparseTableMin(const vector<ll> &a) : table(1,a) {
        int n = a.size();
        for (int half=1,k=1; half<=n/2; half*=2,++k) {
            int count=n-2*half+1;
            table.emplace_back(count);
            for (int i=0;i<count;++i)
                table[k][i] = min(table[k - 1][i],
                                  table[k - 1][i + half]);
        }
    }

    ll query(int l, int r) const { // Requiere 0 <= l <= r < n.
        assert(l<=r); // Longitud positiva: clz(0) no está definido.
        int k = 31-__builtin_clz((unsigned)(r-l+1));
        return min(table[k][l], table[k][r - (1 << k) + 1]);
    }
};
\end{lstlisting}

\begin{lstlisting}[style=cpp]
SparseTableMin rmq(vector<ll>{5,1,4,2,3});
cout << rmq.query(1,3) << '\n'; // 1: incluye los indices 1, 2 y 3.
cout << rmq.query(2,2) << '\n'; // 4: un solo elemento.
\end{lstlisting}

\warning{El arreglo puede ser vacío al construir, pero toda consulta debe
tener índices válidos y rango no vacío. El RMQ original de KACTL usa $[l,r)$;
esta interfaz conserva $[l,r]$. Los valores son \texttt{long long} y no se
modifican al construir. Requiere GCC/Clang para \texttt{\_\_builtin\_clz}.}
